\section{補完空間の非被約性と非純次元性}
\label{sec:ambient-geometry-consequences-v32}

前節までのhull同定は、Poissonスキームの幾何を変形空間へ移す。
ここでは別稿\cite{VKPoisson32}の結果を用い、その変形論的帰結を明記する。
成分と形式局所環の証明は、そこで本稿のhull同定とは独立に与えられる。

\begin{corollary}[非被約な六項補完空間]
\label{cor:ambient-nonreduced-v32}
六項代数 $\Lambda$ のminiversal hullは非被約かつ非純次元であり、完全交叉ではない。
対応する代数的Jacobiスキームは、次元
\[
 35,\ 28,\ 23,\ 22,\ 19,\ 19,\ 19,\ 19
\]
の相異なる既約成分を少なくとも持つ。
最後の十九次元成分の一般点では局所環の長さは六である。
これらの次元はスカラー倍や座標自己同型で商する前の次元である。
\end{corollary}

\begin{proof}
別稿の成分定理と共点三直線の局所模型による。
特に、ある点の完備局所環は滑らかな十八変数を除いて
\[
 \mathbb C[[a,b,u]]/(ab,b^2,a^5u)
\]
となり、$u$ を可逆にすると長さ六の非被約環を得る。
有限型の $\mathbb C$ 代数はexcellentであり、被約局所環の完備化は被約なので、
この局所模型は元のJacobiスキームの非被約性を意味する。
Jacobiイデアルは同次である。非被約軌跡は閉かつスカラー作用で不変であり、
その閉包は原点を含む。従って原点の局所環も非被約である。
完備化の忠実平坦性とhull同定により、六項側のminiversal基底も非被約となる。
異なる次元の円錐成分は全て原点を通るので、この局所環は非純次元である。
斉次素イデアルによる商の原点完備化は整域である。実際、非零な二つの形式元の
最低次斉次項の積は、その斉次整域の中で非零となる。
従って相異なる極小素イデアルとそれらの次元は完備化後にも保たれ、hullも非純次元となる。
正則局所環の完全交叉商は純次元であるから、完全交叉ではない。
\end{proof}

三十八個の独立な二次生成元があることと、局所的な消去の後に五次の式が現れることは両立する。
前者は原点における埋め込みの最小生成次数、後者は非零Poisson点の近傍で
可逆な線形部分を消去した局所方程式の次数である。
また、同じ被約基底上の点集合が一致しても、上の冪零方向はArtin基底上の補完を変える。
本稿の全Artin基底上の節点条件と同様、補完問題ではイデアルの根基への置換は許されない。
